% Independently audited by root and B; source note retained.
\subsection{Acyclic support with at most six excess arrows}

\begin{lemma}[Singleton hom-sets and the support poset]
\label{C:singleton-support-H2}
Let $C$ be a finite connected category with at least two objects,
acyclic support, and exactly one arrow in each nonempty hom-set
between distinct objects.  Let $\Pi$ be its support poset.
The actual functor $q:C\to\Pi$ which collapses every corner
monoid to its identity induces an isomorphism
\[
 q^*:H^2(B\Pi;\mathbb Q)\xrightarrow{\ \cong\ }
      H^2(BC;\mathbb Q).
\]
No assumption on the positive rational homology of the corner
monoids is required.
\end{lemma}
\begin{proof}
The unique cross arrows make $q$ a functor: every composite
between different endpoints is the unique arrow with those
endpoints.  Filter normalized nerve chains by the number $p$
of strict changes of object.  For each chain of $p+1$ distinct
objects, the associated internal bar is the tensor product
of the ordinary bars of their corner monoids, because all
cross coefficients are singleton trivial modules.  Its
degree-zero homology is $\mathbb Q$, so the $q=0$ row is
exactly the simplicial chain complex of $B\Pi$.

Every finite monoid has zero first rational homology.  Hence
the entire internal-degree-one row is zero.  In internal
degree two, the bar for an edge $x<y$ has homology
\[
 H_2(B\operatorname{End}(x);\mathbb Q)
 \oplus H_2(B\operatorname{End}(y);\mathbb Q).
\]
Its first horizontal differential maps these two summands
individually, with the boundary signs, onto the corresponding
corner summands in horizontal degree zero.  Every object
has an incident edge because the support is connected and
has at least two vertices.  Thus $E^2_{0,2}=0$.
The only term of total degree two which remains is
$E^2_{2,0}=H_2(B\Pi;\mathbb Q)$.
No later differential enters it, and its only possible
outgoing differential lands in the already zero term
$E^2_{0,1}$.  Comparison under the actual functor $q$
therefore gives an isomorphism on second rational homology,
and duality gives the displayed cohomology isomorphism.
Naturality makes this a cup-product pullback, not merely
an abstract identification of vector spaces.
\end{proof}

\begin{lemma}[Small cycle rank of a support graph]
\label{C:small-cycle-support}
Let $\Pi$ be a finite connected poset, let $G$ be its
undirected comparability graph, and put
$R=|E(G)|-|V(G)|+1$.
If $R\leq6$, then $H_4(B\Pi;\mathbb Q)=0$.
If $R\leq3$ and $H_1(B\Pi;\mathbb F_2)=0$, then $G$ is
chordal.  In the latter case its maximal-clique decomposition
has the following possibilities after bridges and cut vertices
are separated:
\begin{itemize}
\item at $R=1$, a single triangle;
\item at $R=2$, triangles or two triangles sharing an edge;
\item at $R=3$, at most three triangles glued along edges or
vertices in a clique tree, or one complete four-vertex block.
\end{itemize}
The classifying-space nerve of a category with this support
is the corresponding union of its full clique subcategory
nerves, with their actual full intersection subcategories.
\end{lemma}
\begin{proof}
The order complex of a poset is the clique complex of its
comparability graph: a set of pairwise comparable objects
is a chain.  A four-simplex requires a complete graph on
five vertices, which already has cycle rank six.  If
$R\leq6$ and such a graph occurs, contracting that complete
graph leaves a tree; otherwise the cycle rank would exceed
six.  Consequently this is the only four-simplex, and
there is no five-simplex.  Its nonzero simplicial boundary
shows that fourth homology is zero.  If no such complete
graph occurs, there are no four-simplices at all.

Suppose now $R\leq3$ and $H_1(B\Pi;\mathbb F_2)=0$.
The cycle space of $G$ has dimension $R$ and is spanned
by triangle boundaries.  Choose a basis of at most three
such boundaries.  If an induced cycle of length at least
four existed, each triangle would contain at most one
edge of that cycle: two consecutive edges would give a
forbidden chord.  A sum of at most three basis triangles
therefore cannot contain all the edges of that cycle.
This contradicts the basis property, proving chordality.

The clique-tree description of a chordal graph follows
by removing a simplicial vertex successively.  In the
present small range it can also be read directly from
the cycle budget.  A complete four-vertex block consumes
all three cycles, so its remaining attachments are trees.
Otherwise every nontrivial maximal clique is a triangle;
adjoining a triangle along an edge or a vertex adds one
independent cycle.  There can be at most three triangles.
Two triangles sharing an edge are the diamond graph;
triangles joined only at a vertex belong to different
blocks.  These give exactly the alternatives listed.

Finally the distinct objects used by a nerve simplex
form a chain, hence a clique of $G$.  The simplex therefore
lies in a full maximal-clique subcategory.  Intersections
of these nerves are the nerves of the full intersections.
All inclusions are simplicial subcomplex inclusions, so
their realizations are the indicated homotopy pushouts.
\end{proof}

\begin{theorem}\label{C:acyclic-arrow-excess-six}
Let $C$ be a finite category with $k$ objects, $n$ arrows,
and acyclic support.  If $n\leq2k+5$, then
$BC\not\simeq\mathbb CP^2$.
In particular every category with acyclic support, at
least eight objects, and at most twenty-one arrows is
excluded.
\end{theorem}
\begin{proof}
Suppose to the contrary that $BC\simeq\mathbb CP^2$.
The category is connected.  If $k\leq2$, then $n\leq9$,
contrary to Theorem~\ref{C:CP2-lower16}.  Thus $k\geq3$.
Use the notation
\[
 \begin{split}
 D&=\sum_x(|\operatorname{End}(x)|-1),\\
 P&=\sum_{x<y}(|C(x,y)|-1),\qquad
 R=E-k+1,
 \end{split}
 \qquad D+P+R=n-(2k-1)\leq6,
\]
where the sum is over nonempty cross hom-sets and $E$
counts the edges of the comparability graph.

If $P=0$, Lemma~\ref{C:singleton-support-H2} pulls every
rational degree-two class back from the support poset.
Lemma~\ref{C:small-cycle-support} gives zero fourth
rational homology there.  Its degree-two products, and
therefore their pullbacks, are zero, a contradiction.

If all corner algebras are semisimple over $\mathbb Q$,
the relative bar over their product has no positive
internal Tor.  Since $P\geq1$ gives $R\leq5$, no chain
has five distinct objects.  The relative bar has strict
object degree at most three, so $H_4(BC;\mathbb Q)=0$.
We may suppose henceforth that a nonsemisimple corner
occurs.  An order-one or order-two monoid algebra is
semisimple, so $D\geq2$.

\paragraph{Tree support.}
Suppose $R=0$.  There is no composable pair of strict
arrows, since their composite would create a triangle.
If all corners have zero second rational cohomology,
Corollary~\ref{C:height-one-zero-H2-products} applies.
Otherwise there is exactly one corner $M$ which can have
nonzero second rational cohomology.  Indeed such a corner
has at least five elements by
Theorem~\ref{E:all-six-types}, and two would cost $D\geq8$.
Write $x$ for its object.  If any incident hom-set is a
singleton, its two-object collage is a join and has zero
second cohomology.  Every degree-two class of $BC$ then
restricts to zero on $BM$, and also on every other corner.
Collapsing the corner spaces in the height-one
correspondence pushout pulls all those classes back
from a wedge of suspensions.  Their products are zero.

Thus every edge incident with $x$ has at least two arrows.
If $d$ is its degree in the support tree, then
$d\leq P\leq6-D\leq2$.  When $d=2$, equality forces
$|M|=5$, $D=4$, and two incident hom-sets of size two;
all other corners and hom-sets are singletons.  The
five-element corner with nonzero $H^2$ is the
identity-adjoined $2\times2$ rectangular band.  The two
incident arrows have the same direction, since the
support has height one.  Identifying their two trivial
endpoint objects gives the ordinary four-state act
cone over $M$, or its opposite.  On realization this
is the identification of two points in a connected
CW space, so it induces an isomorphism on cohomology
in degrees at least two.  The act-cone products are
zero by Theorem~\ref{C:two-row-act-cups}; their pullbacks
give zero products on this two-leg star.

When $d=1$, the single incident hom-set has size at most
three.  Size two is covered by
Theorem~\ref{C:two-cross-arrows-universal}.  Size three
forces $|M|=5$, all other corners trivial, and is an
ordinary act cone over the same two-row band monoid;
Theorem~\ref{C:two-row-act-cups} again applies.
In both cases all remaining corners have order at most
two.  The other pieces of the tree have semisimple
diagonal and only one strict crossing, hence zero
rational homology in degrees at least two.  They attach
to the displayed collage or star only at corners of
order at most two.  These intersections are rationally
acyclic.  Mayer--Vietoris therefore makes restriction
in degree four injective and identifies degree-two
classes with their restrictions to the displayed piece.
The square vanishes on that piece and hence on all of
$BC$.  This excludes tree support.

We now have $P,R\geq1$.  Therefore
\[
 2\leq D\leq4,\qquad R\leq3,
\]
and every corner has at most five elements.  All such
corners have $H_3(-;\mathbb Q)=0$ by
Theorem~\ref{E:all-six-types}; those of order at most
four are rationally acyclic.
The actual functor $C\to\Pi$ is onto on fundamental
groups.  Indeed an edge loop in $B\Pi$ lifts by choosing
one actual arrow for each oriented edge, with formal
reversals when necessary.  Since $BC$ is simply
connected, $B\Pi$ is simply connected as well.
Lemma~\ref{C:small-cycle-support} now gives the stated
chordal decompositions.

We repeatedly use two elementary consequences of their
actual nerve pushouts.  A tree of pieces meeting at
rationally acyclic corners has positive rational
homology, in degrees at least two, equal to the direct
sum of the pieces' homology.  Thus each piece inherits
bounded rational homology and second-cohomology
dimension at most one from $BC$.  More generally, if
every intersection has zero third rational homology,
Mayer--Vietoris makes the map from the sum of fourth
homology of the pieces onto fourth homology of the
union surjective.  Equivalently, restriction in fourth
cohomology is injective.  Vanishing of squares on the
pieces then detects vanishing on the union.

\paragraph{One independent cycle.}
Let $R=1$.  The only non-bridge block is a triangle.
If $D\leq3$, all corner intersections are rationally
acyclic.  There is at most one corner of order at least
three, and it has order at most four; all other corners
have order at most two.
For a bridge piece, semisimple corners give zero $H^2$.
Otherwise exact coinvariants at its order-at-most-two
opposite corner give an actual rational cohomology-ring
equivalence to an ordinary act cone over a monoid of
order at most four.  To see the comparison directly,
the edge double bar computes the exact coinvariants
at that corner, and the two vertex bars have the same
rational augmentation.  The double-mapping-cylinder
comparison is consequently a rational equivalence.
The boundedness and rank-one condition inherited by
this act cone permit Lemma~\ref{C:nine-bounded-rank-one},
so its square is zero.

For the triangle, if its larger corner is an endpoint,
that corner has $H^2=0$ and the opposite two-object
block has semisimple diagonal and $H^2=0$.
Lemma~\ref{C:two-block-zero-H2-products} kills its
degree-two products.  If the larger corner is in the
middle, its endpoints are semisimple and its order is
at most four.  The inherited boundedness and
Corollary~\ref{C:semisimple-ends-middle-four} give
zero fourth homology.  A triangle with no larger corner
has semisimple diagonal and zero fourth homology directly.
Fourth-cohomology restriction now excludes a square
on the whole category.

If instead $D=4$, the budget forces $P=1$.  Every bridge
has one or two arrows, so its products vanish by the
join description or Theorem~\ref{C:two-cross-arrows-universal}.
The triangle has at most one doubled hom-set.  If all
its corners have order at most four, at least one
adjacent hom-set is a singleton.  That adjacent block
is a join with zero $H^2$, and the opposite corner
also has zero $H^2$.  The two-block lemma applies.
The only remaining possibility is a five-element corner,
with the other two corners trivial.  If it is the middle
corner, a singleton adjacent hom-set gives the usual
Quillen-A deletion of that middle object, leaving a
graph.  If it is an endpoint and its adjacent hom-set
is a singleton, use the same zero-$H^2$ two-block cut.
If its adjacent hom-set is doubled, the other endpoint
is terminal or initial, so the triangle is contractible.
All corner intersections still have zero $H_3$.
Thus restriction in degree four is injective and again
every global square is zero.

\paragraph{Two independent cycles.}
Let $R=2$.  Now $D\leq3$ and $D+P\leq4$, so all
corner intersections are rationally acyclic.  Bridge
and triangle blocks are handled exactly as in the
preceding $D\leq3$ argument.  The only new block is
a diamond: two triangles sharing an edge $u<v$, with
exclusive vertices $a,b$.  These vertices are
incomparable.  Transitivity forces them to occupy the
same position relative to $u<v$: both below, both
above, or both strictly between.

Call an edge potentially nonzero if its hom-set has
at least two arrows and one endpoint corner has at
least three elements.  Every other edge collage has
$H^2=0$, by the singleton join or semisimple diagonal
argument.  A potentially nonzero edge uses at least
three units of the combined $D+P$ budget.

Suppose first that $a,b<u<v$, the opposite case being
dual.  Cut the diamond into the full blocks $\{a,b\}$
and $\{u,v\}$.  The first has zero $H^2$.  If the shared
edge is not potentially nonzero, the second does too.
Otherwise the shared edge consumes at least three of
the at most four available units.  At least one of
$a,b$ consequently has a trivial corner and singleton
hom-sets to both $u,v$.  It is an initial object in
its full triangle with $u,v$.  Hence every global
degree-two class restricts to zero on the shared-edge
block, as well as on the first block.  The actual
two-block suspension quotient in the proof of
Lemma~\ref{C:two-block-zero-H2-products} pulls back
all these classes, so their products vanish.

Suppose instead $u<a,b<v$.  Consider the two stars
on $\{u,a,b\}$ and $\{a,b,v\}$.  A star has $H^2=0$
if neither edge is potentially nonzero: it is a union
of its two edge collages along a corner with zero
first and second rational homology.  If either star
has zero $H^2$, cut it from the opposite endpoint
corner and apply the two-block lemma.
Avoiding both cuts therefore requires a potentially
nonzero edge in each star.  There is at most one
corner of order at least three, so those edges must
meet at the same exclusive vertex, say $a$.
Their cost is at least two corner excesses and two
parallel-arrow excesses.  Equality with $D+P\leq4$
is forced: the corner at $a$ has order three,
$|C(u,a)|=|C(a,v)|=2$, and every other corner and
cross hom-set in the diamond is a singleton.

The full triangle on $u,b,v$ and its intersection
with the full triangle on $u,a,v$ are contractible;
the intersection is the singleton edge $u\to v$.
The nerve is their actual union, so the inclusion of
the latter triangle induces a homotopy equivalence.
It has semisimple endpoints and middle order three.
Its inherited bounded homology gives zero fourth
homology by Corollary~\ref{C:semisimple-ends-middle-four}.
This closes the diamond.  The corner-intersection
comparison then closes the whole $R=2$ case.

\paragraph{Three independent cycles.}
Finally $R=3$ forces $D=2$ and $P=1$.  There is one
order-three corner, every other corner is trivial,
and exactly one cross hom-set is doubled.  In a
bridge or clique intersection, the only possible
nontrivial corner is this order-three monoid.
Its trivial modules on both sides have projective
dimension at most one, as proved in
Proposition~\ref{C:semisimple-ends-small-middle}.
The edge bar therefore has no internal homology in
degree at least two, and its collage has no rational
homology in degrees at least three.  In particular
every clique intersection has zero third homology.

If the nontrivial cliques are triangles, each has
zero fourth homology.  With the order-three corner
at an endpoint, internal bar degree is at most one
and strict degree at most two.  With that corner
in the middle, at least one adjacent hom-set is a
singleton; the corresponding middle-object deletion
leaves a graph.  With all corners trivial the nerve
is two-dimensional.  Fourth-cohomology restriction
is injective across the clique tree, so this case
has zero fourth homology altogether.

It remains to consider the complete four-vertex
block, whose order is $0<1<2<3$.  A three-object
subchain with trivial corners and at most one doubled
hom-set has zero second homology: its triangle
boundary map is injective, since at least one
adjacent hom-set is a singleton, and the component
on the other adjacent edge distinguishes its
two-simplex labels.
If the order-three corner is outside this block, the
block has only trivial corners and its nerve has
dimension at most three, so its fourth homology is zero.
Otherwise argue as follows.
If the order-three corner is $0$ or $3$, separate
it from the other three objects.  Both full blocks
have zero second cohomology, so the two-block lemma
applies.
If that corner is $1$, first cut into $\{0,1\}$ and
$\{2,3\}$.  This is a zero-$H^2$ cut unless the
doubled hom-set is $C(0,1)$.  In that exception, cut
instead into $\{0\}$ and $\{1,2,3\}$: the second
block has all cross hom-sets singleton and hence
zero $H^2$ by Lemma~\ref{C:singleton-support-H2},
since its support is a simplex.  The case of corner
$2$ is the opposite argument.  Thus every square
vanishes on the four-vertex block.  Its remaining
attachments are bridges with zero third-homology
intersections, so restriction in degree four proves
that every global square vanishes as well.

All possible values of $P,R,D$ have now been covered.
Each contradicts the nonzero rational square of
$\mathbb CP^2$, proving the theorem.  For $k\geq8$
and $n\leq21$ one has $n-(2k-1)\leq6$, giving the
stated consequence.
\end{proof}
